\chapter{Counting and First Probability}

\begin{center}\itshape
$n$ guests leave their hats at the door and get them back at random. The probability that nobody gets their own hat\\
tends to $\frac1e\approx0.368$, whether there are ten guests or ten million.
\end{center}

\section{Counting}

\begin{reminder}[basic counts]
For finite sets with $|E|=n$:
\begin{itemize}
\item maps $F\to E$ with $|F|=k$: $n^k$; injections: $\frac{n!}{(n-k)!}$; bijections $E\to E$ (permutations): $n!$;
\item subsets of size $k$: $\binom nk$; multisets of size $k$ (``stars and bars''): $\binom{n+k-1}k$;
\item \textbf{pigeonhole}: if $m>n$ objects go into $n$ boxes, some box holds at least two (more precisely, at least $\lceil m/n\rceil$).
\end{itemize}
\end{reminder}

\begin{theorem}[inclusion--exclusion]\label{thm:incexc}
For finite sets $A_1,\dots,A_n$, writing $A_I=\bigcap_{i\in I}A_i$,
\[
\left|A_1\cup\dots\cup A_n\right|=\sum_{k=1}^n(-1)^{k+1}\sum_{|I|=k}|A_I|.
\]
\end{theorem}
\begin{proof}
For $x$ in the union, $1-\prod_{i=1}^n\left(1-\mathbf 1_{A_i}(x)\right)=1$; expanding the product,
$\mathbf 1_{\bigcup A_i}=\sum_{\emptyset\ne I}(-1)^{|I|+1}\mathbf 1_{A_I}$, and this also holds (both sides $0$) outside the union. Sum over $x$.
\end{proof}

\begin{proposition}[derangements]\label{prop:derang}
The number $D_n$ of permutations of $\{1,\dots,n\}$ without fixed point is
\[
D_n=n!\sum_{k=0}^n\frac{(-1)^k}{k!},\qquad\text{so}\qquad\frac{D_n}{n!}\to\frac1e.
\]
\end{proposition}
\begin{proof}
Let $A_i$ be the set of permutations fixing $i$. For $|I|=k$, $|A_I|=(n-k)!$. By Theorem~\ref{thm:incexc}, the number of permutations with a fixed point is
$\sum_{k=1}^n(-1)^{k+1}\binom nk(n-k)!=n!\sum_{k=1}^n\frac{(-1)^{k+1}}{k!}$; subtract from $n!$. The limit is the series of $e^{-1}$ (Exercise~\ref{ex:taylorln}).
\end{proof}

\section{Finite probability}

\begin{reminder}[probability on a finite set]
\begin{itemize}
\item A probability on a finite $\Omega$ is $P:\mathcal P(\Omega)\to[0,1]$ with $P(\Omega)=1$ and $P(A\cup B)=P(A)+P(B)$ for disjoint $A,B$. Uniform: $P(A)=|A|/|\Omega|$.
\item Conditional probability $P(A\mid B)=\frac{P(A\cap B)}{P(B)}$; \textbf{Bayes}: $P(B\mid A)=\frac{P(A\mid B)P(B)}{P(A)}$; independence: $P(A\cap B)=P(A)P(B)$.
\item A random variable $X:\Omega\to\R$ has expectation $\mathbb E[X]=\sum_\omega X(\omega)P(\{\omega\})=\sum_xx\,P(X=x)$, and variance $\operatorname{Var}X=\mathbb E[X^2]-\mathbb E[X]^2$.
\item \textbf{Linearity}: $\mathbb E[X+Y]=\mathbb E[X]+\mathbb E[Y]$, with \emph{no independence needed}. For an event $A$, $\mathbb E[\mathbf 1_A]=P(A)$.
\end{itemize}
\end{reminder}

The indicator trick: to compute the expectation of a count, write the count as a sum of indicators and add their probabilities.

\begin{example}
In a uniformly random permutation of $\{1,\dots,n\}$, the expected number of fixed points is $\sum_{i=1}^nP(\sigma(i)=i)=n\cdot\frac1n=1$.
\end{example}

\section{Exercises}

\begin{exercise}[$\star$]\label{ex:stars}
How many solutions in $\N^k$ does $x_1+\dots+x_k=n$ have? And in $(\N^*)^k$?
\end{exercise}

\begin{exercise}[$\star$]\label{ex:birthday}
\emph{The birthday problem.} Show that the probability that $n$ people (with birthdays uniform and independent among 365 days) all have different birthdays is at most
$e^{-n(n-1)/730}$, and deduce that it is less than $\frac12$ for $n=23$. \emph{Use $1-x\le e^{-x}$.}
\end{exercise}

\begin{exercise}[$\star\star$]\label{ex:pigeon}
Choose $n+1$ distinct integers in $\{1,\dots,2n\}$. Show that two of them are coprime, and that one of them divides another.
\end{exercise}

\begin{exercise}[$\star\star$]\label{ex:phi}
Let $n=p_1^{a_1}\cdots p_r^{a_r}$ with distinct primes $p_i$. Using inclusion--exclusion, show that the number $\varphi(n)$ of integers in $\{1,\dots,n\}$ coprime to $n$ is
$n\prod_{i=1}^r\left(1-\frac1{p_i}\right)$.
\end{exercise}

\begin{exercise}[$\star\star$]\label{ex:fixedvar}
Let $X$ be the number of fixed points of a uniformly random permutation of $\{1,\dots,n\}$, $n\ge2$. Show that $\operatorname{Var}X=1$.
\end{exercise}

\begin{exercise}[$\star\star$]\label{ex:cycles}
Let $C$ be the number of cycles of a uniformly random permutation of $\{1,\dots,n\}$.
\begin{enumerate}[label=(\alph*)]
\item Show that the length $L$ of the cycle containing $1$ is uniform on $\{1,\dots,n\}$.
\item Show that $C=\sum_{i=1}^n\frac1{L_i}$, where $L_i$ is the length of the cycle containing $i$, and deduce $\mathbb E[C]=H_n\sim\ln n$.
\end{enumerate}
\end{exercise}

\begin{exercise}[$\star\star\star$]\label{ex:coupon}
\emph{The coupon collector.} Each cereal box contains one of $n$ coupons, uniformly at random and independently. Show that the expected number of boxes needed to collect all $n$
coupons is $nH_n\sim n\ln n$. \emph{You may use that if each trial succeeds with probability $p$ independently, the expected number of trials until the first success is $\frac1p$.}
\end{exercise}

\begin{exercise}[$\star\star\star$]\label{ex:erdos}
\emph{Erd\H{o}s--Szekeres.} Show that every sequence of $n^2+1$ distinct real numbers has a monotone subsequence of length $n+1$.
\emph{Hint: attach to each term the lengths of the longest increasing and decreasing subsequences ending there.}
\end{exercise}

\begin{beauty}
Proposition~\ref{prop:derang} is the first appearance of the Poisson law: the number of fixed points of a random permutation converges to a Poisson variable of
parameter $1$ (Tome~3). The formula for $\varphi$ opens Tome~1, and the indicator trick of Exercises~\ref{ex:fixedvar} to~\ref{ex:coupon} grows into the probabilistic
method of Tome~3, where one proves that an object exists by showing that a random one works with positive probability.
\end{beauty}

\section{Solutions}

\begin{solution}{ex:stars}
A solution is a row of $n$ stars split by $k-1$ bars into $k$ groups (possibly empty), so there are $\binom{n+k-1}{k-1}$. For positive solutions, put $y_i=x_i-1$:
$y_1+\dots+y_k=n-k$ in $\N^k$, giving $\binom{n-1}{k-1}$.
\end{solution}

\begin{solution}{ex:birthday}
The probability is $\prod_{k=0}^{n-1}\left(1-\frac k{365}\right)\le\exp\left(-\sum_{k=0}^{n-1}\frac k{365}\right)=e^{-n(n-1)/730}$.
For $n=23$, $\frac{23\cdot22}{730}=\frac{506}{730}=0.693151\ldots$, just above $\ln2=0.693147\ldots$, so the bound is below $\frac12$.
The exact probability is about $0.4927$.
\end{solution}

\begin{solution}{ex:pigeon}
Split $\{1,\dots,2n\}$ into the $n$ pairs $\{2k-1,2k\}$. Two of the $n+1$ chosen numbers lie in the same pair; they are consecutive, hence coprime.
For divisibility, write each chosen number as $2^am$ with $m$ odd. There are only $n$ odd numbers in $\{1,\dots,2n\}$, so two chosen numbers $2^am$ and $2^bm$ share the same $m$;
the one with the smaller power of $2$ divides the other.
\end{solution}

\begin{solution}{ex:phi}
Let $A_i$ be the set of multiples of $p_i$ in $\{1,\dots,n\}$. For $I\subset\{1,\dots,r\}$, $A_I$ is the set of multiples of $\prod_{i\in I}p_i$, so $|A_I|=n/\prod_{i\in I}p_i$.
An integer is coprime to $n$ if and only if it lies in no $A_i$, so by inclusion--exclusion
\[
\varphi(n)=n-\sum_{\emptyset\ne I}(-1)^{|I|+1}\frac n{\prod_{i\in I}p_i}=n\sum_{I}\prod_{i\in I}\left(-\frac1{p_i}\right)=n\prod_{i=1}^r\left(1-\frac1{p_i}\right).
\]
\end{solution}

\begin{solution}{ex:fixedvar}
Write $X=\sum_i\mathbf 1_{\sigma(i)=i}$. Then $X^2=\sum_i\mathbf 1_{\sigma(i)=i}+\sum_{i\ne j}\mathbf 1_{\sigma(i)=i,\,\sigma(j)=j}$. A permutation fixing $i\ne j$ is a permutation of the other $n-2$ points,
so $P(\sigma(i)=i,\sigma(j)=j)=\frac{(n-2)!}{n!}$. Hence $\mathbb E[X^2]=1+n(n-1)\frac{(n-2)!}{n!}=2$ and $\operatorname{Var}X=2-1^2=1$.
\end{solution}

\begin{solution}{ex:cycles}
(a) A permutation whose cycle through $1$ has length $\ell$ is given by the ordered list $\sigma(1),\sigma^2(1),\dots,\sigma^{\ell-1}(1)$ of distinct elements other than $1$
($\frac{(n-1)!}{(n-\ell)!}$ choices) and an arbitrary permutation of the $n-\ell$ remaining points ($(n-\ell)!$ choices): $(n-1)!$ permutations in all, for each $\ell$. So $P(L=\ell)=\frac{(n-1)!}{n!}=\frac1n$.

(b) Each cycle of length $\ell$ contributes $\ell\cdot\frac1\ell=1$ to $\sum_i\frac1{L_i}$, so the sum counts cycles. By symmetry every $L_i$ has the law of $L$, and by linearity
$\mathbb E[C]=n\,\mathbb E\left[\frac1L\right]=n\cdot\frac1n\sum_{\ell=1}^n\frac1\ell=H_n$, which is $\sim\ln n$ by Proposition~\ref{prop:gamma}.
\end{solution}

\begin{solution}{ex:coupon}
Let $T_k$ be the number of boxes bought while holding exactly $k$ distinct coupons, $0\le k\le n-1$. During that phase each box brings a new coupon with probability $\frac{n-k}n$,
independently, so $\mathbb E[T_k]=\frac n{n-k}$. The total is $T=\sum_kT_k$, and by linearity
$\mathbb E[T]=\sum_{k=0}^{n-1}\frac n{n-k}=n\sum_{j=1}^n\frac1j=nH_n$.
\end{solution}

\begin{solution}{ex:erdos}
For each term $a_i$, let $u_i$ (resp.\ $d_i$) be the length of the longest increasing (resp.\ decreasing) subsequence ending at $a_i$. If $i<j$ then $(u_i,d_i)\ne(u_j,d_j)$:
if $a_i<a_j$, appending $a_j$ gives $u_j\ge u_i+1$; if $a_i>a_j$, then $d_j\ge d_i+1$. If all the $u_i,d_i$ were at most $n$, the $n^2+1$ distinct pairs would lie in
$\{1,\dots,n\}^2$, which has $n^2$ elements: impossible by pigeonhole. So some $u_i$ or $d_i$ is at least $n+1$.
\end{solution}
